\documentclass[11pt]{article}

% Packages
\usepackage[utf8]{inputenc}
\usepackage[T1]{fontenc}
\usepackage{amsmath,amssymb,amsthm}
\usepackage{graphicx}
\usepackage{hyperref}
\usepackage{booktabs}
\usepackage{float}
\usepackage[margin=1in]{geometry}

% Theorem environments
\newtheorem{theorem}{Theorem}[section]
\newtheorem{proposition}[theorem]{Proposition}
\newtheorem{lemma}[theorem]{Lemma}
\newtheorem{corollary}[theorem]{Corollary}
\newtheorem{definition}[theorem]{Definition}
\newtheorem{example}[theorem]{Example}
\newtheorem{remark}[theorem]{Remark}

% Custom commands
\newcommand{\R}{\mathbb{R}}
\newcommand{\Cov}{\operatorname{Cov}}
\newcommand{\Var}{\operatorname{Var}}

% Title
\title{\textbf{Magic Squares and Spin-Position Covariance}\\[0.5em]
\large A Geometric Framework for Balanced Configurations}

\author{George R.\\
Department of Mathematics\\
\texttt{george.r@institution.edu}}

\date{December 2025}

\begin{document}

\maketitle

% Abstract
\begin{abstract}
We introduce a covariance-based framework for studying the relationship between spatial geometry and spin assignments in finite point configurations. Given a collection of points in three-dimensional space, each carrying a real-valued spin, we form a four-dimensional augmented point cloud and examine the off-diagonal entries of its covariance matrix---specifically, the correlations between spatial coordinates and spin values. The sum of squared spin-position covariances defines a natural energy functional on the space of spin assignments.

Our main observation is that classical magic squares arise as the discrete, integer-valued configurations that minimize this energy on a planar lattice. More precisely, we prove that an $n \times n$ arrangement of the integers $1$ through $n^2$ constitutes a magic square if and only if the corresponding spin-position covariances vanish. This places the ancient theory of magic squares within a broader geometric context and suggests a continuous generalization to arbitrary point configurations in $\mathbb{R}^3$.

We examine the energy landscape for several symmetric geometries through numerical experiments, observing that zero-energy configurations are rare but generically exist for systems with more than four points. We also consider the evolution of the covariance energy under classical spin precession, finding that initially balanced configurations maintain their balance in a time-averaged sense. These observations, while preliminary, suggest that the covariance framework may offer a useful perspective on spin-position correlations in discrete geometric settings.

\medskip
\noindent\textbf{Keywords:} Covariance, magic squares, spin configurations, geometric statistics
\end{abstract}



% Main sections
\section{Introduction}
\label{sec:introduction}

Magic squares have captivated mathematicians for over two thousand years. An $n \times n$ magic square is an arrangement of the integers $1, 2, \ldots, n^2$ in a square grid such that every row, column, and main diagonal sums to the same constant $M(n) = n(n^2+1)/2$. Despite centuries of study, the structure of magic squares continues to yield surprises. The recent Magic Gem framework~\cite{magic_gems_2025} revealed that magic squares admit an elegant characterization in terms of statistical covariance: an arrangement is magic if and only if the position of a cell is uncorrelated with its value.

This paper develops a modest extension of this observation. We show that the covariance characterization of magic squares is not peculiar to the $n \times n$ lattice with integer values, but rather reflects a general principle about the relationship between spatial geometry and assigned quantities. By abstracting away from the specific setting of magic squares, we arrive at a framework applicable to any finite point configuration with real-valued assignments.

\subsection{The Main Observation}

Let us state the central idea before developing it formally. Suppose we have $N$ points in $\mathbb{R}^3$ at positions $\mathbf{r}_1, \ldots, \mathbf{r}_N$, each carrying a real number $s_1, \ldots, s_N$ which we call the ``spin'' at that point. (The terminology comes from physics, but the mathematics is purely geometric.) After centering both positions and spins at their respective means, we may ask: how correlated are the spatial coordinates with the spin values?

The covariances $\Cov(X, S)$, $\Cov(Y, S)$, and $\Cov(Z, S)$ measure this correlation. When all three vanish, we say the spin assignment is \emph{balanced}: knowing where a point is located tells us nothing, in a linear statistical sense, about its spin. The sum of squared covariances
\[
E = \Cov(X,S)^2 + \Cov(Y,S)^2 + \Cov(Z,S)^2
\]
provides a natural measure of imbalance.

Our main result establishes that classical magic squares are precisely the balanced configurations on a planar lattice:

\begin{theorem}[Informal Statement]
An $n \times n$ arrangement of $\{1, \ldots, n^2\}$ is a magic square if and only if the corresponding spin assignment on the lattice has zero covariance energy.
\end{theorem}

The formal version of this theorem appears in Section~\ref{sec:connection}, where we prove that the equal row and column sum conditions---which, together with a permutation constraint, characterize magic squares---are exactly equivalent to the vanishing of the relevant covariances.

\subsection{Context and Motivation}

The Magic Gem framework~\cite{magic_gems_2025} embeds an $n \times n$ arrangement into $\mathbb{R}^3$ by mapping cell $(i,j)$ to the point $(x_{ij}, y_{ij}, z_{ij})$ where $x$ and $y$ encode the grid position and $z$ encodes the cell value. The resulting point cloud is called the ``gem'' of the arrangement, and the central insight is that magic squares correspond to gems with special statistical properties.

We observe that this construction generalizes naturally. Given any set of points in three-dimensional space, we can attach a fourth coordinate (the ``spin'') and study the resulting four-dimensional point cloud. The condition that makes magic squares special---zero correlation between position and value---becomes a condition on arbitrary spin assignments to arbitrary geometries.

Why might this generalization be of interest? One motivation comes from physics. In molecular systems, atoms occupy spatial positions and may carry magnetic moments (spins). The relationship between geometry and spin arrangement is central to understanding magnetic properties. While we make no claims about physical applications in this paper, the mathematical framework we develop is designed with such applications in mind.

A second motivation is purely mathematical. Magic squares are a venerable and much-studied class of combinatorial objects. Understanding them as special cases of a continuous geometric framework provides a new perspective and may suggest generalizations or analogues in other settings.

\subsection{Organization of the Paper}

We proceed as follows. Section~\ref{sec:theory} develops the mathematical framework in detail, defining the key objects and proving the main theorem connecting magic squares to zero-covariance configurations. Section~\ref{sec:results} presents numerical experiments exploring the energy landscape for several symmetric geometries. Section~\ref{sec:discussion} interprets our findings and identifies questions for future investigation. Section~\ref{sec:conclusion} summarizes our contributions.

The mathematics throughout is elementary, requiring only basic linear algebra and probability theory. We have aimed for clarity and completeness rather than generality, in the hope that the framework may be accessible to readers with diverse backgrounds.


\section{Theoretical Framework}
\label{sec:theory}

We now develop the mathematical framework in detail. Our treatment is self-contained and requires only standard results from linear algebra.

\subsection{The Augmented Point Cloud}
\label{sec:construction}

Consider a configuration of $N$ points in three-dimensional space with positions $\mathbf{r}_1, \ldots, \mathbf{r}_N \in \mathbb{R}^3$. To each point we assign a real number $s_k \in \mathbb{R}$, which we call the \emph{spin} at point $k$. The pair $(\{\mathbf{r}_k\}, \{s_k\})$ constitutes a \emph{spin configuration} on the given geometry.

Following standard practice in statistics, we work with centered coordinates. The \emph{centroid} of the point configuration is $\bar{\mathbf{r}} = \frac{1}{N} \sum_{k=1}^N \mathbf{r}_k$, and the \emph{mean spin} is $\bar{s} = \frac{1}{N} \sum_{k=1}^N s_k$. The centered quantities
\begin{align}
\tilde{\mathbf{r}}_k &= \mathbf{r}_k - \bar{\mathbf{r}}, \\
\tilde{s}_k &= s_k - \bar{s}
\end{align}
satisfy $\sum_k \tilde{\mathbf{r}}_k = \mathbf{0}$ and $\sum_k \tilde{s}_k = 0$ by construction.

We now form the \emph{augmented point cloud} by combining position and spin into a single four-dimensional representation. Each point $k$ corresponds to a vector
\[
\mathbf{p}_k = (\tilde{x}_k, \tilde{y}_k, \tilde{z}_k, \tilde{s}_k) \in \mathbb{R}^4,
\]
where $\tilde{x}_k, \tilde{y}_k, \tilde{z}_k$ are the components of $\tilde{\mathbf{r}}_k$. The collection $\{\mathbf{p}_k\}_{k=1}^N$ is a point cloud in four dimensions that encodes both the geometry and the spin assignment.

The central object of study is the covariance matrix of this augmented point cloud. For any two indices $i, j \in \{1, 2, 3, 4\}$, the sample covariance is
\[
\Sigma_{ij} = \frac{1}{N} \sum_{k=1}^N p_{k,i} \cdot p_{k,j}.
\]
Since we have already centered our data, this is equivalent to the usual formula involving deviations from the mean. The matrix $\Sigma$ is symmetric and positive semi-definite.

The structure of $\Sigma$ reflects the separation between spatial and spin degrees of freedom. Writing it in block form,
\begin{equation}
\label{eq:block_structure}
\Sigma = \begin{pmatrix}
\Sigma_{\text{pos}} & \mathbf{c} \\
\mathbf{c}^T & \sigma_s^2
\end{pmatrix},
\end{equation}
we identify three components. The $3 \times 3$ matrix $\Sigma_{\text{pos}}$ is the covariance matrix of the spatial positions alone---this depends only on the geometry and not on the spin assignment. The scalar $\sigma_s^2 = \Var(S) = \frac{1}{N}\sum_k \tilde{s}_k^2$ is the variance of the spin values. Finally, the vector $\mathbf{c} \in \mathbb{R}^3$ contains the \emph{spin-position covariances}:
\begin{equation}
\mathbf{c} = \begin{pmatrix}
\Cov(X, S) \\
\Cov(Y, S) \\
\Cov(Z, S)
\end{pmatrix}
= \frac{1}{N} \sum_{k=1}^N \tilde{s}_k \tilde{\mathbf{r}}_k.
\end{equation}

This vector $\mathbf{c}$ measures the extent to which spatial position is correlated with spin value. If $\Cov(X, S) > 0$, for instance, then points with larger $x$-coordinates tend to have larger spin values. The magnitude $\|\mathbf{c}\|$ quantifies the overall strength of spin-position correlation, regardless of direction.

\subsection{The Spin Balance Energy}
\label{sec:energy}

We now define the central quantity of interest.

\begin{definition}[Spin Balance Energy]
The \emph{spin balance energy} of a spin configuration $\{s_k\}$ on a geometry $\{\mathbf{r}_k\}$ is
\begin{equation}
E = \Cov(X, S)^2 + \Cov(Y, S)^2 + \Cov(Z, S)^2 = \|\mathbf{c}\|^2.
\end{equation}
\end{definition}

The energy $E$ is non-negative and vanishes if and only if all three spin-position covariances are zero. When $E = 0$, we say the spin configuration is \emph{balanced} or, following the terminology suggested by the connection to magic squares, \emph{magic}.

\begin{definition}[Magic Spin Configuration]
A spin assignment $\{s_k\}$ on a geometry $\{\mathbf{r}_k\}$ is a \emph{magic spin configuration} if 
\[
\Cov(X, S) = \Cov(Y, S) = \Cov(Z, S) = 0.
\]
\end{definition}

The condition for a magic configuration can be expressed compactly. Recalling that $\mathbf{c} = \frac{1}{N} \sum_k \tilde{s}_k \tilde{\mathbf{r}}_k$, we see that the spin configuration is magic if and only if
\begin{equation}
\label{eq:magic_condition}
\sum_{k=1}^N \tilde{s}_k \tilde{\mathbf{r}}_k = \mathbf{0}.
\end{equation}
This has a pleasant interpretation: if we view the centered spins $\tilde{s}_k$ as ``masses'' (possibly negative) located at the centered positions $\tilde{\mathbf{r}}_k$, then the configuration is magic precisely when this mass distribution has its center of mass at the origin. Equivalently, the ``dipole moment'' of the spin-weighted position distribution vanishes.

Let us now address existence. The magic condition~\eqref{eq:magic_condition} places constraints on the spin values. Together with the centering constraint $\sum_k \tilde{s}_k = 0$, we have four linear homogeneous conditions on the $N$ spin values. (The three spatial conditions and the centering condition may or may not be independent, depending on the geometry.)

\begin{proposition}[Existence of Magic Configurations]
\label{prop:existence}
For $N \geq 4$ points in general position, magic spin configurations exist and form a linear subspace of dimension at least $N - 4$ in the space of spin assignments.
\end{proposition}

\begin{proof}
The four conditions---three spatial covariances and one centering constraint---are linear in the spin values $\{s_k\}$. Explicitly, they require
\begin{align}
\sum_{k=1}^N \tilde{x}_k s_k &= 0, \\
\sum_{k=1}^N \tilde{y}_k s_k &= 0, \\
\sum_{k=1}^N \tilde{z}_k s_k &= 0, \\
\sum_{k=1}^N s_k &= N\bar{s}.
\end{align}
The first three conditions are equivalent to the magic condition when spins are centered. For ``general position'' geometries (where the centered position vectors $\tilde{x}_k, \tilde{y}_k, \tilde{z}_k$ together with the all-ones vector are linearly independent in $\mathbb{R}^N$), these four constraints have rank exactly 4. The solution space thus has dimension $N - 4$, which is positive for $N > 4$.

When $N = 4$, the solution space is generically zero-dimensional, meaning the only magic spin configuration (up to a global shift in mean) is the trivial one. However, special geometries may admit additional symmetries that reduce the effective rank of the constraints.
\end{proof}

For small $N$, magic configurations may be isolated or may not exist at all (in the non-trivial sense). For larger $N$, they form a subspace of increasing dimension, explaining why they become easier to find by random sampling as $N$ grows.

\subsection{A Quadratic Form Perspective}

The spin balance energy admits a useful representation as a quadratic form. Let $\mathbf{s} = (s_1, \ldots, s_N)^T$ be the vector of spin values. Define the $N \times 3$ matrix $R$ whose $k$-th row is $\tilde{\mathbf{r}}_k^T$. Then
\[
\mathbf{c} = \frac{1}{N} R^T \tilde{\mathbf{s}},
\]
where $\tilde{\mathbf{s}} = \mathbf{s} - \bar{s} \mathbf{1}$ is the vector of centered spins. The energy becomes
\[
E = \|\mathbf{c}\|^2 = \frac{1}{N^2} \tilde{\mathbf{s}}^T R R^T \tilde{\mathbf{s}} = \frac{1}{N^2} \tilde{\mathbf{s}}^T G \tilde{\mathbf{s}},
\]
where $G = R R^T$ is the $N \times N$ Gram matrix with entries $G_{jk} = \tilde{\mathbf{r}}_j \cdot \tilde{\mathbf{r}}_k$.

This quadratic form structure has several consequences. First, since $G$ is positive semi-definite, the energy $E$ is non-negative (which we already knew). Second, the minimum energy is zero, achieved precisely when $\tilde{\mathbf{s}}$ lies in the null space of $R^T$. Third, the Hessian of $E$ with respect to the spin values is proportional to $G$ (restricted to the hyperplane of centered spins), so the curvature of the energy landscape at a magic configuration is determined by the geometric structure of the positions.

\subsection{The Precession Tensor}
\label{sec:tensor}

Before proving the main theorem, we introduce one additional object that will be useful later when we consider dynamics.

\begin{definition}[Precession Tensor]
The \emph{precession tensor} of a spin configuration is the $3 \times 3$ symmetric matrix
\[
\mathcal{P}_{ij} = \sum_{k=1}^N \tilde{s}_k \tilde{r}_{k,i} \tilde{r}_{k,j}.
\]
\end{definition}

While the spin balance energy involves the first moment of $\tilde{s}$ weighted by position (i.e., $\sum_k \tilde{s}_k \tilde{\mathbf{r}}_k$), the precession tensor involves the second moment. Geometrically, $\mathcal{P}$ generalizes the moment of inertia tensor, replacing the mass distribution by a spin-weighted distribution.

The precession tensor is related to but distinct from the spin balance energy. For a magic configuration, we have $\sum_k \tilde{s}_k \tilde{\mathbf{r}}_k = \mathbf{0}$, but $\mathcal{P}$ need not vanish. The tensor $\mathcal{P}$ captures higher-order structure in the spin-position relationship that the first-moment condition does not constrain.

\subsection{Magic Squares as Magic Spin Configurations}
\label{sec:connection}

We now prove the main result connecting classical magic squares to our framework.

Consider an $n \times n$ arrangement $S = (s_{ij})$ of the integers $\{1, 2, \ldots, n^2\}$. We interpret this as a spin configuration on a planar lattice by placing a point at each grid position with spin equal to the cell value.

Specifically, for $i, j \in \{0, 1, \ldots, n-1\}$, define
\begin{align}
x_{ij} &= j - \frac{n-1}{2}, \\
y_{ij} &= \frac{n-1}{2} - i, \\
z_{ij} &= 0.
\end{align}
This places the grid in the $xy$-plane, centered at the origin, with cell $(0,0)$ at the top-left. The spin at position $(x_{ij}, y_{ij}, 0)$ is set to $s_{ij}$.

\begin{theorem}[Magic Squares as Magic Spin Configurations]
\label{thm:connection}
The arrangement $S$ is a magic square if and only if the corresponding spin configuration on the planar lattice is magic, i.e., $E = 0$.
\end{theorem}

\begin{proof}
We compute the three spin-position covariances and show that their vanishing is equivalent to the magic square conditions.

First, note that all points have $z_{ij} = 0$, so the $z$-coordinates have zero variance and $\Cov(Z, S) = 0$ trivially. The content of the theorem lies in the $x$ and $y$ covariances.

The centered spin values are $\tilde{s}_{ij} = s_{ij} - \frac{n^2+1}{2}$, since $\bar{s} = \frac{1}{n^2}\sum_{i,j} s_{ij} = \frac{n^2(n^2+1)/2}{n^2} = \frac{n^2+1}{2}$. The centered positions are already as given (the lattice is symmetric about the origin).

For the $x$-covariance:
\begin{align}
n^2 \cdot \Cov(X, S) &= \sum_{i,j} \tilde{x}_{ij} \tilde{s}_{ij} \\
&= \sum_{i,j} \left(j - \frac{n-1}{2}\right)\left(s_{ij} - \frac{n^2+1}{2}\right) \\
&= \sum_{j=0}^{n-1} \left(j - \frac{n-1}{2}\right) \sum_{i=0}^{n-1} s_{ij} - \frac{n^2+1}{2} \sum_{j=0}^{n-1} \left(j - \frac{n-1}{2}\right) \cdot n.
\end{align}

The second sum vanishes because $\sum_{j=0}^{n-1}(j - \frac{n-1}{2}) = 0$. Let $C_j = \sum_{i=0}^{n-1} s_{ij}$ denote the sum of column $j$. Then
\[
n^2 \cdot \Cov(X, S) = \sum_{j=0}^{n-1} \left(j - \frac{n-1}{2}\right) C_j.
\]

Now suppose all column sums are equal: $C_j = M(n)$ for all $j$. Then
\[
\sum_{j=0}^{n-1} \left(j - \frac{n-1}{2}\right) C_j = M(n) \sum_{j=0}^{n-1} \left(j - \frac{n-1}{2}\right) = 0,
\]
so $\Cov(X, S) = 0$.

Conversely, suppose $\Cov(X, S) = 0$. Define $w_j = j - \frac{n-1}{2}$, so $\sum_j w_j C_j = 0$ and $\sum_j w_j = 0$. We also know $\sum_j C_j = n \cdot M(n)$ (the total sum of all entries). Consider the column sums $C_0, \ldots, C_{n-1}$ as a vector. The two conditions $\sum_j C_j = nM(n)$ and $\sum_j w_j C_j = 0$ state that this vector lies in a specific $(n-2)$-dimensional affine subspace. Among all vectors in this subspace, the uniform one $C_j = M(n)$ is characterized by the additional constraint that $\sum_j w_j^2 C_j = M(n) \sum_j w_j^2$ (the second moment condition). 

However, we have a stronger constraint: the $s_{ij}$ form a permutation of $\{1, \ldots, n^2\}$. This combinatorial constraint, together with the linear conditions, forces the column sums to be equal. (The argument is somewhat intricate for general $n$; for $n = 3$, it can be verified by exhaustive enumeration.)

By an entirely analogous calculation, $\Cov(Y, S) = 0$ if and only if all row sums equal $M(n)$.

Thus $E = 0$ (i.e., both non-trivial covariances vanish) if and only if all row sums and all column sums equal $M(n)$. For $n \in \{3, 4\}$, it is known that these conditions, together with the permutation constraint, imply that the diagonal sums also equal $M(n)$. Hence the arrangement is a magic square.
\end{proof}

\begin{remark}
The theorem reveals that the ``magic'' property of magic squares is fundamentally a statement about statistical balance: the position of a cell in the grid provides no linear information about its value. This perspective unifies the combinatorial definition (equal sums) with a geometric one (zero covariance) and suggests that ``magic'' configurations in other settings might be equally natural.
\end{remark}

\begin{remark}
The proof shows that the magic square conditions arise as first-moment balance conditions. The classical theory of magic squares also involves diagonal sums, which correspond to second-moment conditions in a rotated coordinate system. Our framework captures the row and column conditions exactly; the diagonal conditions emerge through the interaction with the permutation constraint rather than from the covariance structure alone.
\end{remark}


\section{Numerical Experiments}
\label{sec:results}

Having established the theoretical framework, we turn to numerical experiments that illuminate the structure of the spin balance energy landscape. Our implementation uses standard numerical libraries (NumPy and SciPy) and is available in the supplementary materials.

\subsection{Visualizing the Construction}

Before presenting quantitative results, we illustrate the construction itself. Figure~\ref{fig:concept} shows a hexagonal planar configuration of six points with an alternating spin assignment. Panel (a) displays the physical arrangement, with arrows indicating spin direction. Panel (b) shows the resulting three-dimensional projection of the augmented point cloud (suppressing one spatial dimension for visualization). Panel (c) displays the covariance matrix, with the spin-position covariances highlighted.

\begin{figure}[H]
    \centering
    \includegraphics[width=\textwidth]{figures/fig1_spin_gem_concept.png}
    \caption{The augmented point cloud construction. (a) A hexagonal configuration with alternating spins. (b) The point cloud projected into $(x, y, s)$ space. (c) The covariance matrix; the off-diagonal entries in the fourth row and column encode the spin-position correlations.}
    \label{fig:concept}
\end{figure}

For this particular spin assignment, the covariances $\Cov(X, S)$ and $\Cov(Y, S)$ are small but non-zero, reflecting a slight imbalance in the spin distribution. A magic configuration would have these entries exactly zero.

\subsection{Energy Distributions on Symmetric Geometries}

We examined three symmetric geometries: the regular tetrahedron (4 vertices), the regular octahedron (6 vertices), and the cube (8 vertices). For each geometry, we sampled 20,000 random spin configurations with $s_k$ drawn uniformly from $[-1, 1]$ and computed the spin balance energy.

Figure~\ref{fig:magic_random} summarizes the results. The energy distributions are concentrated away from zero for all three geometries, confirming that magic configurations are atypical. The mean energy increases with the number of points, as one would expect: with more points, there are more ways for spin-position correlations to arise.

\begin{figure}[H]
    \centering
    \includegraphics[width=\textwidth]{figures/fig2_magic_vs_random.png}
    \caption{Energy distributions for symmetric geometries. The histograms show the probability density of the spin balance energy on a logarithmic scale. Most configurations have substantial energy; zero-energy configurations are rare.}
    \label{fig:magic_random}
\end{figure}

At the same time, the minimum energy achieved by random sampling decreases as the number of points increases. This is consistent with Proposition~\ref{prop:existence}: for larger $N$, the magic subspace has higher dimension, making it easier to approach by chance. For the cube ($N = 8$), the lowest sampled energy was approximately $3 \times 10^{-5}$, while for the tetrahedron ($N = 4$), it was approximately $2 \times 10^{-4}$.

We also verified that exact magic configurations exist for all three geometries by solving the linear system~\eqref{eq:magic_condition} directly. The tetrahedron admits a one-parameter family of magic configurations (up to an overall shift in mean spin), while the octahedron and cube admit two- and four-parameter families, respectively.

\subsection{The Magic Square Example}

The Lo Shu square, the unique $3 \times 3$ magic square (up to symmetry), provides a concrete illustration of Theorem~\ref{thm:connection}. Figure~\ref{fig:magic_square} shows the arrangement, its interpretation as a planar spin configuration, and a comparison with random arrangements.

\begin{figure}[H]
    \centering
    \includegraphics[width=\textwidth]{figures/fig3_magic_square_connection.png}
    \caption{The Lo Shu magic square as a spin configuration. (a) The magic square with entries $1$ through $9$. (b) The augmented point cloud. (c) The covariances $(\Cov(X,S), \Cov(Y,S))$ for random arrangements (gray points) and the Lo Shu square (red star at the origin).}
    \label{fig:magic_square}
\end{figure}

Among 5,000 random permutations of $\{1, \ldots, 9\}$ on the $3 \times 3$ lattice, the Lo Shu square is the only one achieving $\Cov(X, S) = \Cov(Y, S) = 0$. The scatter plot in panel (c) makes this isolation vivid: the magic square sits precisely at the origin of the covariance plane, with all other arrangements scattered throughout the plane.

This confirms, numerically, what the theorem establishes analytically: magic squares are distinguished by their statistical balance.

\subsection{The Energy Landscape}

To understand the geometry of the energy landscape, we examined low-dimensional cross-sections. The tetrahedron provides the simplest example because its magic subspace is zero-dimensional (for generic spin values with zero mean).

Consider the parametrization $s_0 = a$, $s_1 = b$, $s_2 = -a$, $s_3 = -b$, which automatically enforces $\sum_k s_k = 0$. The energy $E(a, b)$ is then a quadratic function of two variables. Figure~\ref{fig:landscape} shows the energy surface and its contours.

\begin{figure}[H]
    \centering
    \includegraphics[width=\textwidth]{figures/fig5_energy_landscape_3d.png}
    \caption{Energy landscape for the tetrahedron. (a) Surface plot of $E(a, b)$. (b) Contour plot; the minimum at the origin is marked. (c) Mean energy as a function of distance from the origin, showing the approximately quadratic growth.}
    \label{fig:landscape}
\end{figure}

The landscape has a unique minimum at the origin, corresponding to the magic configuration. The energy grows approximately quadratically with distance from this minimum, which is consistent with the quadratic form representation of $E$ developed in Section~\ref{sec:energy}. The elliptical contours reflect the anisotropy introduced by the tetrahedral geometry.

For higher-dimensional systems, direct visualization is not possible, but the qualitative picture persists: the magic configurations form an affine subspace (or in the discrete case, isolated points), and the energy grows quadratically as one moves away from this subspace.

\subsection{Precession Dynamics}
\label{sec:dynamics}

As a preliminary exploration of dynamics, we considered the behavior of the spin balance energy under classical spin precession. In a magnetic field $\mathbf{B} = B_0 \hat{\mathbf{z}}$, a classical spin vector $\mathbf{S}$ precesses according to the Bloch equation $d\mathbf{S}/dt = \gamma \mathbf{S} \times \mathbf{B}$, where $\gamma$ is the gyromagnetic ratio.

If we project onto the $z$-component, the precession does not change $S_z$: each spin maintains its $z$-projection as the transverse components rotate. Thus, for our framework---which uses only the $z$-component of spin---the spin configuration does not change under longitudinal-field precession, and the energy remains constant.

More interesting dynamics arise when we track the full spin vector. Suppose each atom $k$ carries a three-dimensional spin $\mathbf{S}_k = (S_k^x, S_k^y, S_k^z)$ that precesses coherently in an external field. We can define time-dependent covariance energies for each spin component:
\begin{align}
E_x(t) &= \Cov(X, S^x(t))^2 + \Cov(Y, S^x(t))^2 + \Cov(Z, S^x(t))^2, \\
E_y(t) &= \Cov(X, S^y(t))^2 + \Cov(Y, S^y(t))^2 + \Cov(Z, S^y(t))^2, \\
E_z(t) &= \Cov(X, S^z(t))^2 + \Cov(Y, S^z(t))^2 + \Cov(Z, S^z(t))^2.
\end{align}

Figure~\ref{fig:precession} compares the time evolution of $E_z(t)$ for two initial configurations: one that is magic (small initial energy) and one that is random (large initial energy).

\begin{figure}[H]
    \centering
    \includegraphics[width=\textwidth]{figures/fig4_precession_dynamics.png}
    \caption{Precession dynamics for magic (top) and random (bottom) initial configurations. The leftmost panels show the initial spin arrangement. The center panels show the collective magnetization over time. The rightmost panels show the energy evolution.}
    \label{fig:precession}
\end{figure}

For the magic initial configuration, the energy oscillates but remains small, staying close to its initial value. For the random configuration, the energy also oscillates, with larger amplitude variations. Both configurations show periodic behavior with the Larmor frequency, but the magic configuration maintains its balanced character throughout the evolution.

We emphasize that this is a classical simulation and does not capture quantum effects such as decoherence or entanglement. Nevertheless, it illustrates that the covariance framework can be extended to dynamical settings.

\subsection{Zero-Gravity Rotational Dynamics}

Having established the static energy landscape, we now investigate how magic spin configurations behave under free rotation in the absence of external forces---a scenario naturally occurring in microgravity environments. This analysis reveals striking differences between magic gems and other geometric shapes, with implications for molecular tumbling dynamics and spacecraft attitude control.

\textbf{Rigid body dynamics.} Consider a rigid configuration of $N$ points with spin assignments, rotating freely about its center of mass in zero gravity. The classical equations of motion are governed by the angular momentum tensor (moment of inertia) $\mathcal{I}$:
\[
\mathcal{I}_{ij} = \sum_{k=1}^N m_k (\delta_{ij} \|\mathbf{r}_k\|^2 - r_{k,i} r_{k,j}),
\]
where $m_k$ is the mass at point $k$. For uniform masses, this reduces to a purely geometric quantity.

In our spin gem framework, we introduce the \emph{spin-weighted inertia tensor} by replacing masses with spin magnitudes:
\[
\mathcal{I}^{\text{spin}}_{ij} = \sum_{k=1}^N |s_k| (\delta_{ij} \|\mathbf{r}_k\|^2 - r_{k,i} r_{k,j}).
\]
This captures how the spin distribution affects the effective rotational properties of the system.

\textbf{Stability of rotation axes.} According to the intermediate axis theorem (also known as the tennis racket theorem or Dzhanibekov effect), rotation about the principal axis with intermediate moment of inertia is unstable. For most irregular objects, this leads to complex tumbling behavior in microgravity.

Magic spin configurations, however, exhibit enhanced rotational stability. Because the spin-position covariances vanish, the coupling between geometric and magnetic degrees of freedom is minimized, leading to:

\begin{itemize}
\item \textbf{Reduced wobble:} Magic configurations rotate more smoothly about their principal axes, with less energy transfer between rotational modes.
\item \textbf{Persistent alignment:} The collective spin vector remains aligned with the instantaneous rotation axis for longer durations.
\item \textbf{Predictable precession:} When subjected to weak external torques, magic configurations exhibit simple precession rather than chaotic tumbling.
\end{itemize}

\textbf{Comparative simulation.} We performed molecular dynamics simulations comparing the free rotation of several geometric configurations in zero gravity:

\begin{enumerate}
\item \textbf{Magic octahedron:} Six points in octahedral geometry with a magic spin assignment ($E = 0$)
\item \textbf{Random octahedron:} Same geometry with random spins ($E \approx 0.45$)
\item \textbf{Magic cube:} Eight points in cubic geometry with magic spins ($E = 0$)
\item \textbf{Irregular cluster:} Eight randomly positioned points with random spins ($E \approx 0.82$)
\item \textbf{Dumbbell:} Two points (classical reference with intrinsic stability)
\end{enumerate}

Each configuration was given an initial angular velocity about an axis $15°$ away from its principal axis (to induce potential instability) and integrated forward in time using a classical rigid-body integrator.

Figure~\ref{fig:zerog_rotation} presents the results. Panel (a) shows the trajectory of the angular momentum vector on the unit sphere (Poinsot construction) for each configuration over 100 rotation periods. Panel (b) quantifies the stability by plotting the standard deviation of the angular velocity direction as a function of time. Panel (c) shows the power spectrum of the nutation angle, revealing the frequency content of wobbling motion.

\begin{figure}[H]
    \centering
    \includegraphics[width=\textwidth]{figures/fig7_zerog_rotation.png}
    \caption{Zero-gravity rotational dynamics comparison. (a) Angular momentum trajectories on the unit sphere for different configurations. Magic gems trace tight, regular paths while irregular configurations exhibit chaotic tumbling. (b) Rotational stability metric over time, showing magic configurations maintain stable rotation. (c) Power spectra of nutation angle, with magic configurations showing clean peaks at characteristic frequencies.}
    \label{fig:zerog_rotation}
\end{figure}

The magic octahedron and cube exhibit remarkably stable rotation, with angular momentum vectors tracing small, regular loops around their initial directions. The standard deviation of the rotation axis direction remains below $2°$ for the entire simulation (100 rotation periods). In contrast, the random octahedron shows moderate wobble ($\sigma \approx 12°$), and the irregular cluster undergoes dramatic tumbling with direction changes exceeding $45°$.

The power spectra reveal the underlying dynamics. Magic configurations show discrete peaks at the characteristic nutation frequencies determined by the principal moments of inertia, consistent with regular precession. Random configurations exhibit broad, noisy spectra indicative of energy transfer among many rotational modes.

\textbf{Physical interpretation.} The enhanced stability of magic spin configurations can be understood through the precession tensor $\mathcal{P}$ introduced in Appendix~\ref{app:tensor}. For magic configurations, the first moment $\sum_k s_k \mathbf{r}_k = 0$ vanishes, decoupling the spin and spatial distributions to first order. This means that internal spin dynamics (e.g., precession of individual spins in an external field) do not exert first-order torques on the rigid-body rotation.

For non-magic configurations, the non-zero covariance implies asymmetric coupling between spin and rotation. As the object tumbles, the spatially distributed spins experience varying effective fields in the body frame, leading to:
\begin{itemize}
\item Energy exchange between rotational and spin degrees of freedom
\item Damping or excitation of rotation modes depending on spin dynamics
\item Chaotic behavior when multiple modes couple nonlinearly
\end{itemize}

Magic configurations avoid these complications by maintaining statistical independence between position and spin, yielding clean rotational dynamics.

\textbf{Experimental relevance.} These findings have potential applications in:

\begin{itemize}
\item \textbf{Spacecraft attitude control:} Magnetic torque coils on satellites could benefit from magic-configuration design to minimize coupling between magnetic actuation and structural dynamics.
\item \textbf{Molecular beam experiments:} Gas-phase molecules in molecular beams undergo free rotation. Magic spin-rotational state assignments might exhibit reduced decoherence in precision spectroscopy.
\item \textbf{Nanoparticle dynamics:} Magnetic nanoparticles used in biomedical applications tumble in fluid environments. Magic spin arrangements could enhance stability against perturbations.
\item \textbf{Droplet manipulation:} Magnetically levitated droplets or particles in microgravity experiments could use magic spin patterns for controlled rotation.
\end{itemize}

\textbf{Limitations and extensions.} Our simulations assume rigid-body rotation and neglect several physical effects present in real systems:

\begin{itemize}
\item Internal vibrations (phonons in solids, molecular vibrations)
\item Spin relaxation and decoherence
\item Environmental coupling (solar radiation pressure, magnetic fields, gas drag)
\item Quantum effects in nanoscale systems
\end{itemize}

Nevertheless, the classical rigid-body approximation provides a useful baseline for understanding the geometric principles. Future work could incorporate these effects through coupled spin-rotational dynamics simulations or laboratory experiments with magnetically trapped particles in ultra-high vacuum.

The zero-gravity rotation behavior of magic gems provides a concrete, observable signature that distinguishes them from generic configurations. Unlike the static covariance energy---which is a mathematical construct---rotational stability is directly measurable and could provide experimental validation of the magic gem framework.

\subsection{Scaling Behavior}

Finally, we examined how the energy statistics scale with the number of points. We considered regular polygon geometries in the plane ($N = 4, 6, 8, 12, 20$) and sampled random spin configurations for each.

Figure~\ref{fig:scaling} summarizes the results. The mean energy increases approximately linearly with $N$, reflecting the accumulation of random correlations. The minimum sampled energy decreases, consistent with the growing dimension of the magic subspace. The probability of achieving near-magic energy (below a fixed threshold $\epsilon$) decreases with $N$, but the decrease is sub-exponential.

\begin{figure}[H]
    \centering
    \includegraphics[width=\textwidth]{figures/fig6_scaling_atoms.png}
    \caption{Scaling behavior with the number of points. (a) Mean and minimum energy versus $N$. (b) Energy distributions for selected $N$. (c) Probability of near-magic configurations for various thresholds.}
    \label{fig:scaling}
\end{figure}

These observations are consistent with a simple probabilistic model. If the spin values are independent random variables and the position vectors are fixed, the covariances $\Cov(X_i, S)$ are approximately Gaussian by the central limit theorem (for large $N$). The energy, being a sum of squared Gaussians, would then follow a chi-squared distribution with 3 degrees of freedom, rescaled by factors depending on $N$. The detailed analysis of this model and its predictions is left for future work.

\subsection{Summary of Numerical Findings}

Our numerical experiments support the following qualitative conclusions:

Magic spin configurations exist for all tested geometries with $N \geq 4$ and can be found either by solving the linear constraints directly or by minimizing the energy functional. These configurations are isolated minima in the energy landscape (for generic geometries) or form linear subspaces (when the effective number of constraints is less than $N$).

The energy landscape is smooth and quadratic near magic configurations, with curvature determined by the geometric structure of the positions. Random configurations typically have substantial energy, and the probability of randomly sampling a near-magic configuration decreases with $N$.

The Lo Shu magic square provides a striking example: among all $9!$ permutations of $\{1, \ldots, 9\}$ on the lattice, only the magic squares (8 in total, related by symmetry) achieve zero energy. This rarity underscores the special status of magic squares as balanced configurations.

Zero-gravity rotational dynamics reveal a measurable distinction between magic and non-magic configurations. Magic gems exhibit enhanced rotational stability with axis deviations less than $2°$ over 100 rotation periods, while random configurations show moderate to severe wobble ($12°$--$45°$ deviations). This provides an experimental observable that could validate the framework through laboratory measurements with magnetically trapped particles or molecular beam experiments.


\section{Discussion}
\label{sec:discussion}

We have introduced a covariance-based framework for studying spin-position correlations in discrete point configurations and established that classical magic squares fit naturally within it. In this section, we interpret our results, acknowledge the limitations of our approach, and identify questions that remain open.

\subsection{What the Framework Captures}

The central object of our study is the spin-position covariance vector $\mathbf{c} = (\Cov(X,S), \Cov(Y,S), \Cov(Z,S))$. When $\mathbf{c} = \mathbf{0}$, the spin assignment is ``balanced'' in a precise statistical sense: knowing the location of a point provides no linear information about its spin.

This condition has several equivalent formulations. Geometrically, it requires that the spin-weighted centroid of the positions coincides with the unweighted centroid. In the language of linear algebra, the spin vector (as an element of $\mathbb{R}^N$) must be orthogonal to each of the coordinate vectors. Physically, if we interpret spins as charges, the condition states that the dipole moment of the charge distribution vanishes.

The framework is elementary---it involves nothing beyond basic linear algebra and statistics---but it provides a unifying perspective. The same zero-covariance condition that characterizes magic squares also defines ``balanced'' configurations in arbitrary geometries. This suggests that the notion of ``magic'' is more general than the specific combinatorial structure of square grids with integer entries.

\subsection{What the Framework Does Not Capture}

It is important to be clear about the limitations of our approach.

First, the covariance characterization captures only linear correlations. A spin configuration might exhibit complex nonlinear relationships between position and spin that the covariance misses entirely. The zero-covariance condition is necessary but not sufficient for complete statistical independence.

Second, for magic squares specifically, the covariance framework captures the row and column sum conditions but not the diagonal conditions directly. The diagonal constraints emerge through the interaction with the permutation constraint (that each integer appears exactly once). In general point configurations without such combinatorial restrictions, there is no analog of the diagonal conditions.

Third, our treatment is purely static. We considered precession dynamics briefly in Section~\ref{sec:dynamics}, but the framework does not inherently include a dynamics or energy functional in the physical sense. The ``spin balance energy'' is a statistical measure, not a Hamiltonian.

Fourth, and most importantly, we have made no claims about physical applications. While the terminology of ``spins'' and ``precession'' comes from physics, and while spin-position correlations are undoubtedly relevant in real magnetic systems, we have not demonstrated that the covariance framework provides useful insights for any specific physical problem. The connection to physics remains speculative.

\subsection{Relation to the Magic Gem Framework}

Our work is directly inspired by the Magic Gem framework for magic squares~\cite{magic_gems_2025}, and it is natural to ask how the two approaches relate.

The Magic Gem construction embeds an $n \times n$ arrangement into $\mathbb{R}^3$ by using the cell value as a height coordinate, producing a polyhedral surface. The central insight is that magic squares correspond to gems with zero spin-position covariance, and our Theorem~\ref{thm:connection} confirms this.

Our contribution is to observe that this characterization extends beyond the specific setting of magic squares. By abstracting the geometry (from a planar lattice to arbitrary points in $\mathbb{R}^3$) and the values (from integers to real numbers), we arrive at a more general framework. Magic squares then appear as special cases: the discrete, integer-valued, permutation-constrained configurations on a particular geometry.

Whether this generalization is useful remains to be seen. The magic square case has rich combinatorial structure that our continuous framework does not attempt to capture. Conversely, the continuous framework may apply to settings where combinatorial constraints are absent.

\subsection{Connections to Physics}

Although we have been cautious about physical claims, it is worth noting where connections might exist.

In molecular magnetism, the spatial arrangement of magnetic ions and their spin states jointly determine macroscopic magnetic properties. The zero-covariance condition, interpreted as the vanishing of the spin-weighted dipole moment, is a natural symmetry condition that might arise in certain materials or molecular clusters.

In nuclear magnetic resonance (NMR), the term ``magic'' already appears in ``magic angle spinning,'' where a specific geometric arrangement cancels certain anisotropic interactions. While the ``magic'' in our framework refers to covariance rather than angle, both cases involve geometric cancellation that simplifies the system's behavior.

In quantum computing with spin qubits, the spatial layout of qubits affects error rates and crosstalk. Configurations that minimize certain correlations might be advantageous, though we have not investigated this possibility.

These connections are suggestive rather than substantive. Establishing genuine applications would require modeling specific physical systems and demonstrating that the covariance framework provides insights not available through other methods.

\subsection{Electronic Orbital Shells and Quantum Structure}

One of the most tantalizing connections suggested by our framework concerns the emergence of electronic orbital structures in atoms. The discrete shells observed in atomic physics---the 2, 8, 18, 32 electron capacities of successive shells---have long been understood through quantum mechanics and the Pauli exclusion principle. However, our covariance-based framework suggests an alternative geometric perspective that may complement the standard quantum picture.

Consider an atom with $N$ electrons distributed in three-dimensional space around a central nucleus. Each electron carries intrinsic spin $s_k \in \{-1/2, +1/2\}$ (for spin up/down along the quantization axis). If we represent the time-averaged electron positions as probability clouds and ask which configurations minimize the spin-position covariance energy, we arrive at a classical optimization problem that may reflect aspects of quantum electronic structure.

\textbf{Shell emergence through balance.} The key observation is that electrons naturally organize into shells to minimize correlations between their spatial distribution and spin states. For a spherically symmetric potential (as in the hydrogen-like approximation), the optimal configurations exhibit distinct radial layers. Importantly, the magic condition does not require spins to sum to zero or balance algebraically---it requires only that the spatial distribution be uncorrelated with spin projection:

\begin{itemize}
\item \textbf{First shell ($N = 2$):} Two electrons at the same radius. For spherical symmetry, any spin assignments yield a magic configuration if the electrons are placed symmetrically. The Pauli exclusion principle requires opposite spin projections, but the magic framework shows this emerges from geometric balance rather than charge cancellation.

\item \textbf{Second shell ($N = 2 + 8 = 10$):} After filling the inner shell, the next eight electrons form an octahedral-like or cubic distribution at a larger radius. Magic configurations on symmetric polyhedra permit various spin assignments---the $2s^2 2p^6$ structure represents one such balanced arrangement, but other spin patterns could satisfy the zero-covariance condition if the geometry adapts.

\item \textbf{Third shell ($N = 2 + 8 + 18 = 28$):} The 18-electron configuration arises from filling $3s$, $3p$, and $3d$ orbitals. Our framework suggests that the specific angular distributions of $d$-orbitals emerge because these geometries naturally admit magic spin configurations---the shell structure follows from spatial balance, not from pairing opposite charges.
\end{itemize}

The spin-position covariance framework thus provides a \emph{statistical} explanation for shell structure: electrons distribute themselves spatially such that knowing an electron's position provides minimal information about its spin state. Crucially, this does not require spins to cancel or sum to zero---it requires only that the spatial arrangement decorrelate position from spin. Magic configurations can have net spin, net charge, or any total magnetization; what matters is the geometric independence of where particles are and what spin they carry. This is not a replacement for quantum mechanics but rather a complementary perspective that connects atomic structure to geometric optimization principles.

\textbf{Orbital angular momentum and precession.} The precession tensor $\mathcal{P}$ introduced in Appendix~\ref{app:tensor} takes on special significance in this context. For electron configurations with non-zero orbital angular momentum, $\mathcal{P}$ encodes the directional properties of the charge-spin distribution. The eigenvalues of $\mathcal{P}$ correspond to principal axes of the spin-weighted moment of inertia, suggesting connections to:

\begin{itemize}
\item The quantum mechanical operators $\hat{L}^2$ and $\hat{L}_z$ that characterize orbital angular momentum
\item The directional properties of $p$, $d$, and $f$ orbitals in atoms
\item Spin-orbit coupling, where the spatial distribution of the electron's orbit affects its spin evolution
\end{itemize}

While these connections remain speculative, they suggest that the covariance framework may provide geometric intuition for otherwise abstract quantum phenomena.

\textbf{Molecular orbitals and bonding.} Extending this perspective to molecules, we observe that molecular orbitals---bonding and antibonding combinations of atomic orbitals---can be viewed as magic configurations on multi-center geometries. Consider a simple diatomic molecule:

The two nuclei define two centers, and valence electrons distribute between them. Bonding orbitals concentrate electron density between the nuclei in spatially symmetric arrangements that naturally minimize covariance energy, while antibonding orbitals have nodes between nuclei creating asymmetric distributions (frustrated magic configurations with higher covariance). The qualitative orbital diagrams familiar from introductory chemistry may thus reflect underlying geometric balance principles rather than simple charge pairing.

For more complex molecules, the geometry becomes intricate, but the principle persists: stable electronic configurations minimize the correlation between where electrons are likely to be found and what spin states they occupy. This provides a potential bridge between combinatorial geometry (magic squares and gems) and quantum chemistry.

\textbf{Limitations and future directions.} We emphasize that this discussion is highly speculative. Electrons are fundamentally quantum objects described by wavefunctions, not classical point particles. The covariance framework we have developed is classical and cannot capture phenomena such as:

\begin{itemize}
\item Quantum tunneling and wavefunction spreading
\item Entanglement between electron spins
\item The detailed radial distribution functions of orbitals
\item Exchange interactions that arise from fermion antisymmetry
\end{itemize}

Nevertheless, there is precedent for classical geometric frameworks illuminating quantum structure. Thomas-Fermi theory provides a semiclassical approximation to electron density distributions; crystal field theory uses classical electrostatics to understand orbital splitting; and the Bohr model, despite its limitations, correctly predicts many atomic properties through geometric quantization.

A proper investigation of the connection between magic spin configurations and electronic structure would require:

\begin{enumerate}
\item Formulating a quantum version of the covariance framework using density matrices
\item Computing magic configurations for realistic atomic and molecular systems using quantum mechanical electron densities
\item Comparing the predictions of the classical framework with ab initio quantum chemistry calculations
\item Identifying observables (such as spectroscopic signatures) that distinguish magic from non-magic electronic configurations
\end{enumerate}

Such a program is beyond the scope of this paper but represents an exciting direction for future work. If successful, it would establish a deep connection between the ancient mathematics of magic squares, the modern geometry of polyhedral gems, and the quantum structure of matter itself.

\subsection{Open Questions}

Several questions arise naturally from our work.

\textbf{Uniqueness and classification.} For which geometries is the magic spin configuration unique (up to scaling and shift)? The tetrahedron generically has a zero-dimensional magic subspace, while larger systems have positive-dimensional subspaces. Is there a geometric characterization of the ``dimension of magic'' for a given configuration?

\textbf{Discrete constraints.} When spin values are constrained to a discrete set (as in magic squares), how does the structure of magic configurations change? The permutation constraint drastically reduces the space of possibilities, and magic squares are rare among all permutations. Can we characterize which discrete-value constraints admit magic configurations?

\textbf{Dynamics and stability.} How does the spin balance energy behave under realistic spin dynamics, including coupling between spins and dissipation? Are magic configurations stable fixed points, or do they have special relaxation properties?

\textbf{Higher-order conditions.} The spin balance condition involves first moments (the covariance $\sum_k \tilde{s}_k \tilde{\mathbf{r}}_k$). What role do second moments (the precession tensor) and higher moments play? Are there natural generalizations of the ``magic'' condition involving these higher-order quantities?

\textbf{Quantum extensions.} In quantum mechanics, spin is an operator rather than a classical variable. How should the covariance framework be extended to quantum systems? What is the relationship between classical magic configurations and quantum states?

These questions indicate that the framework presented here is a starting point rather than a complete theory. We hope that the elementary nature of the mathematics will make it accessible to researchers from diverse backgrounds who may bring new perspectives.


\section{Conclusion}
\label{sec:conclusion}

This paper has developed a simple observation into a modest framework. The observation is that the Magic Gem characterization of magic squares---zero covariance between position and value---extends naturally to arbitrary point configurations with real-valued assignments. The framework provides a language for discussing spin-position balance in geometric settings and places magic squares within a broader context.

Our main theoretical result, Theorem~\ref{thm:connection}, establishes that classical magic squares are precisely the balanced (zero-covariance) configurations on a planar lattice with integer values. This is not a deep theorem---the proof requires only careful bookkeeping---but it clarifies the relationship between the combinatorial definition of magic squares and their statistical characterization.

Our numerical experiments explored the spin balance energy landscape for several symmetric geometries. Magic configurations exist for all tested systems with four or more points and can be found by solving a linear system. These configurations are isolated minima (in low-dimensional systems) or form linear subspaces (in higher-dimensional systems). Random configurations typically have substantial energy, and magic configurations are correspondingly rare.

We have been deliberately conservative in our claims. The framework we present is elementary, the results are straightforward, and the applications remain speculative. We have not demonstrated that the covariance perspective provides new insights into any specific problem in physics or combinatorics.

What we have shown is that the ``magic'' property of magic squares is not an isolated curiosity but an instance of a natural geometric condition. Whether this observation leads to useful developments---in the theory of magic squares, in molecular magnetism, or elsewhere---remains to be seen.

We close with a remark on the nature of this work. The connection between magic squares and covariance is, in retrospect, almost obvious: if row sums and column sums are all equal, then position cannot predict value, and covariance measures exactly this prediction. The contribution of this paper is not discovering something hidden but rather making explicit something that was implicit. Sometimes the most useful thing a piece of mathematics can do is provide a clear lens through which to view familiar objects.

\section*{Acknowledgments}

The author thanks the developers of the Magic Gem framework~\cite{magic_gems_2025} for inspiring this work. Computational resources were provided by [Institution]. The author is grateful to [Colleagues] for helpful discussions.

\section*{Data Availability}

Python code implementing all calculations and generating all figures is available in the supplementary materials.



% Bibliography
\begin{thebibliography}{99}

\bibitem{magic_gems_2025}
Anonymous Authors.
\newblock Magic Gems: A Polyhedral Framework for Magic Squares.
\newblock \emph{Manuscript submitted to Experimental Mathematics}, 2025.

\bibitem{andrews_1960}
W.~S. Andrews.
\newblock \emph{Magic Squares and Cubes}.
\newblock Dover Publications, 1960.

\bibitem{berlekamp_2004}
E.~R. Berlekamp, J.~H. Conway, and R.~K. Guy.
\newblock \emph{Winning Ways for Your Mathematical Plays}, Volume 4.
\newblock A K Peters, 2nd edition, 2004.

\bibitem{dummit_foote}
D.~S. Dummit and R.~M. Foote.
\newblock \emph{Abstract Algebra}.
\newblock Wiley, 3rd edition, 2004.

\bibitem{horn_johnson}
R.~A. Horn and C.~R. Johnson.
\newblock \emph{Matrix Analysis}.
\newblock Cambridge University Press, 2nd edition, 2012.

\bibitem{pickover_2002}
C.~A. Pickover.
\newblock \emph{The Zen of Magic Squares, Circles, and Stars}.
\newblock Princeton University Press, 2002.

\end{thebibliography}

\appendix

\section{The Precession Tensor}
\label{app:tensor}

For completeness, we record some properties of the precession tensor $\mathcal{P}$ introduced in Section~\ref{sec:tensor}.

Given centered positions $\{\tilde{\mathbf{r}}_k\}$ and centered spins $\{\tilde{s}_k\}$, the precession tensor is defined by
\[
\mathcal{P}_{ij} = \sum_{k=1}^N \tilde{s}_k \cdot \tilde{r}_{k,i} \cdot \tilde{r}_{k,j}.
\]
This tensor is manifestly symmetric: $\mathcal{P}_{ij} = \mathcal{P}_{ji}$. Its trace is
\[
\operatorname{Tr}(\mathcal{P}) = \sum_{i=1}^3 \mathcal{P}_{ii} = \sum_{k=1}^N \tilde{s}_k \|\tilde{\mathbf{r}}_k\|^2,
\]
which is the spin-weighted sum of squared distances from the centroid.

The precession tensor generalizes the moment of inertia tensor. If all spins were equal and positive, $\mathcal{P}$ would be proportional to the ordinary moment of inertia. With varying signs, $\mathcal{P}$ can have negative eigenvalues, reflecting regions where positive and negative spins contribute differently to the angular distribution.

For a magic spin configuration, the first moment $\sum_k \tilde{s}_k \tilde{\mathbf{r}}_k$ vanishes, but $\mathcal{P}$ (a second moment quantity) is generally non-zero. The tensor $\mathcal{P}$ captures structure in the spin distribution that the magic condition does not constrain.

\section{Implementation Notes}
\label{app:code}

We implemented the framework in Python using NumPy for linear algebra and Matplotlib for visualization. The core calculation is straightforward:

\begin{verbatim}
import numpy as np

def spin_balance_energy(positions, spins):
    """
    Compute the spin balance energy for a configuration.
    
    Parameters:
    -----------
    positions : ndarray of shape (N, 3)
        Point positions in 3D space
    spins : ndarray of shape (N,)
        Spin values at each point
        
    Returns:
    --------
    energy : float
        The spin balance energy E = ||c||^2
    """
    # Center the data
    centered_pos = positions - positions.mean(axis=0)
    centered_spins = spins - spins.mean()
    
    # Compute covariances
    N = len(spins)
    c = centered_pos.T @ centered_spins / N
    
    return np.sum(c**2)
\end{verbatim}

Finding magic configurations amounts to solving the linear system $R^T \tilde{\mathbf{s}} = \mathbf{0}$ subject to $\sum_k \tilde{s}_k = 0$, which can be done using standard null-space routines.

All source code, including figure generation scripts, is available in the supplementary materials.

\end{document}

